\renewcommand{\abstractname}{摘\quad 要}
\renewcommand{\keywordsname}{关键词}
\renewcommand{\subjclassname}{2020 年数学主题分类}
\renewcommand{\contentsname}{目\quad 录}
\renewcommand{\proofname}{证明}
\renewcommand{\refname}{参考文献}
\makeatletter
\def\emailaddrname{{\itshape 通讯邮箱}}
\makeatother

\crefname{section}{第~}{第~}
\crefname{subsection}{第~}{第~}
\crefname{subsubsection}{第~}{第~}
\crefname{equation}{式}{式}
\crefname{figure}{图}{图}
\crefname{table}{表}{表}
\crefname{assumption}{假设}{假设}
\crefname{remark}{注}{注}
\crefname{problem}{问题}{问题}
\crefname{proposition}{命题}{命题}
\crefname{theorem}{定理}{定理}
\crefname{lemma}{引理}{引理}
\crefname{corollary}{推论}{推论}
\crefname{definition}{定义}{定义}
\crefname{conjecture}{猜想}{猜想}
\crefname{example}{例}{例}
\crefname{part}{第~}{第~}
\AtBeginDocument{%
% Chinese formatting for single, multiple, and ranged section references.
% Keep the suffix only once, e.g. “第 1、2 节” rather than “第 1 节、第 2 节”.
\crefformat{section}{第~#2#1#3~节}
\Crefformat{section}{第~#2#1#3~节}
\crefmultiformat{section}{第~#2#1#3}{、#2#1#3~节}{、#2#1#3}{、#2#1#3~节}
\Crefmultiformat{section}{第~#2#1#3}{、#2#1#3~节}{、#2#1#3}{、#2#1#3~节}
\crefrangeformat{section}{第~#3#1#4~至~#5#2#6~节}
\Crefrangeformat{section}{第~#3#1#4~至~#5#2#6~节}

\crefformat{subsection}{第~#2#1#3~节}
\Crefformat{subsection}{第~#2#1#3~节}
\crefmultiformat{subsection}{第~#2#1#3}{、#2#1#3~节}{、#2#1#3}{、#2#1#3~节}
\Crefmultiformat{subsection}{第~#2#1#3}{、#2#1#3~节}{、#2#1#3}{、#2#1#3~节}
\crefrangeformat{subsection}{第~#3#1#4~至~#5#2#6~节}
\Crefrangeformat{subsection}{第~#3#1#4~至~#5#2#6~节}

\crefformat{subsubsection}{第~#2#1#3~节}
\Crefformat{subsubsection}{第~#2#1#3~节}
\crefmultiformat{subsubsection}{第~#2#1#3}{、#2#1#3~节}{、#2#1#3}{、#2#1#3~节}
\Crefmultiformat{subsubsection}{第~#2#1#3}{、#2#1#3~节}{、#2#1#3}{、#2#1#3~节}
\crefrangeformat{subsubsection}{第~#3#1#4~至~#5#2#6~节}
\Crefrangeformat{subsubsection}{第~#3#1#4~至~#5#2#6~节}

  \renewcommand{\crefpairconjunction}{、}%
  \renewcommand{\crefmiddleconjunction}{、}%
  \renewcommand{\creflastconjunction}{、}%
  \renewcommand{\crefrangeconjunction}{至}%
  \renewcommand{\subjclassname}{2020 年数学主题分类}%
  \renewcommand{\appendixname}{附录}%
  \renewcommand{\tablename}{表}%
  \renewcommand{\figurename}{图}%
}
